\documentclass[../../main.tex]{subfiles}% 注意这里的文件路径不能用 ./main.tex ,否则用latexmk编译子文件会报错
\graphicspath{{\subfix{./image/}}{\subfix{../../image/}}} % 指定图片引用路径,后续可以直接使用图片文件名
% 如果图片存放在上级目录,则使用 ../../image/ 作为备用路径

% 例如:
% \begin{figure}[H]
% \centering
% \includegraphics[width=0.8\textwidth]{图.png}
% \caption{}
% \label{figure:图}
% \end{figure}
% 注意:上述\label{}一定要放在\caption{}之后,否则引用图片序号会只会显示??.

\begin{document}

\section{三重积分}

\begin{theorem}[Gauss公式]\label{theorem:Gauss公式--数分}
设 $\Omega \subset \mathbb{R}^3$ 是一个有界闭区域，其边界 $\partial\Omega$ 为分片光滑的定向曲面，并取外法向 $\bm{n}=(n_x,n_y,n_z)$。设向量场
\[
\bm{F}(x,y,z) = \bigl( P(x,y,z),\, Q(x,y,z),\, R(x,y,z) \bigr)
\]
在包含 $\overline{\Omega}$ 的某个开集上具有连续的一阶偏导数,即 $\bm{F}\in C^1(\overline{\Omega})$。则成立
\[
\iiint_{\Omega}{\left( \frac{\partial P}{\partial x}+\frac{\partial Q}{\partial y}+\frac{\partial R}{\partial z} \right) \,\mathrm{d}V}=\iiint_{\Omega}{\nabla \cdot \boldsymbol{F}\,\mathrm{d}V}=\iint_{\partial \Omega}{\boldsymbol{F}\cdot \boldsymbol{n}\,\mathrm{d}S}=\iint_{\partial \Omega}{\left( P\cdot n_x+Q\cdot n_y+R\cdot n_z \right) \mathrm{d}S}.
\]
其中 $\nabla \cdot \bm{F}= \frac{\partial P}{\partial x} + \frac{\partial Q}{\partial y} + \frac{\partial R}{\partial z}$ 为 $\bm{F}$ 的散度，$\mathrm{d}V$ 为体积元，$\mathrm{d}S$ 为面积元。上式也可以写成微分形式:
\begin{align*}
\iiint_{\Omega}{\left( \frac{\partial P}{\partial x}+\frac{\partial Q}{\partial y}+\frac{\partial R}{\partial z} \right) \mathrm{d}x\mathrm{d}y\mathrm{d}z}=\iint_{\partial \Omega}{\bigl( P\mathrm{d}y\mathrm{d}z+Q\mathrm{d}z\mathrm{d}x+R\mathrm{d}x\mathrm{d}y \bigr)},
\end{align*}
其中$\mathrm{d}y\mathrm{d}z=n_x\mathrm{d}S,  \mathrm{d}z\mathrm{d}x=n_y\mathrm{d}S,  \mathrm{d}x\mathrm{d}y=n_z\mathrm{d}S$.
\end{theorem}
\begin{remark}
若 $\partial\Omega$ 由多个连通分支组成，则各分支均取外法向。
\end{remark}

\begin{theorem}[Stokes公式]\label{theorem:Stokes(斯托克斯)公式--数分}
设 $S \subset \mathbb{R}^3$ 是一个分片光滑的有向曲面，其边界 $\partial S$ 是一条或有限条分段光滑的闭曲线，且方向与 $S$ 的定向满足右手法则（即当拇指指向 $S$ 的法向时，四指方向为 $\partial S$ 的正向）。设向量场
\[
\bm{F}(x,y,z) = \bigl( P(x,y,z),\, Q(x,y,z),\, R(x,y,z) \bigr)
\]
在包含 $S$ 的某个开集上具有连续的一阶偏导数,即 $\bm{F}\in C^1(\overline{S})$。则成立
\begin{align}\label{eqIOUIOFWJIOFJIOJWIRJIOWJ93482984028094844}
\iint_{S} (\nabla \times \bm{F}) \cdot \bm{n} \, \mathrm{d}S
= \oint_{\partial S} \bm{F}\cdot \mathrm{d}\bm{r},
\end{align}
其中 $\nabla \times \bm{F}$ 为 $\bm{F}$ 的旋度：
\[
\nabla \times \bm{F}= 
\begin{vmatrix}
\mathbf{i} & \mathbf{j} & \mathbf{k} \\
\partial_x & \partial_y & \partial_z \\
P & Q & R
\end{vmatrix}
= \left( \frac{\partial R}{\partial y} - \frac{\partial Q}{\partial z},\,
\frac{\partial P}{\partial z} - \frac{\partial R}{\partial x},\,
\frac{\partial Q}{\partial x} - \frac{\partial P}{\partial y} \right),
\]
$\bm{n}=(n_x,n_y,n_z)$ 是 $S$ 的单位法向量（与 $S$ 的定向一致），$\mathrm{d}\bm{r}$ 为切向线元。\eqref{eqIOUIOFWJIOFJIOJWIRJIOWJ93482984028094844}式展开得
\begin{align*}
\iint_S{\left[ \left( \frac{\partial R}{\partial y}-\frac{\partial Q}{\partial z} \right) n_x+\left( \frac{\partial P}{\partial z}-\frac{\partial R}{\partial x} \right) n_y+\left( \frac{\partial Q}{\partial x}-\frac{\partial P}{\partial y} \right) n_z \right] \mathrm{d}S}=\oint_{\partial S}{P\mathrm{d}x+Q\mathrm{d}y+R\mathrm{d}z}.
\end{align*}
上式也可以写成微分形式:
\begin{align*}
\iint_S{\left( \frac{\partial R}{\partial y}-\frac{\partial Q}{\partial z} \right) \mathrm{d}y\mathrm{d}z+\left( \frac{\partial P}{\partial z}-\frac{\partial R}{\partial x} \right) \mathrm{d}x\mathrm{d}z+\left( \frac{\partial Q}{\partial x}-\frac{\partial P}{\partial y} \right) \mathrm{d}x\mathrm{d}y}=\oint_{\partial S}{P\mathrm{d}x+Q\mathrm{d}y+R\mathrm{d}z},
\end{align*}
其中$\mathrm{d}y\mathrm{d}z=n_x\mathrm{d}S,  \mathrm{d}z\mathrm{d}x=n_y\mathrm{d}S,  \mathrm{d}x\mathrm{d}y=n_z\mathrm{d}S$.
\end{theorem}
\begin{remark}
当 $S$ 是平面区域且 $\bm{n}$ 为常数时，Stokes公式退化为二维Green公式.
\end{remark}
















\end{document}